\documentclass[11pt, a4paper]{article}

% ─── Paquetes ────────────────────────────────────────────────────────────────
\usepackage[utf8]{inputenc}
\usepackage[T1]{fontenc}
\usepackage[spanish]{babel}
\usepackage{amsmath, amssymb, amsthm}
\usepackage{geometry}
\usepackage{xcolor}
\usepackage{mdframed}
\usepackage{titlesec}
\usepackage{enumitem}
\usepackage{booktabs}
\usepackage{array}
\usepackage{hyperref}
\usepackage{ifthen}

\geometry{top=2.2cm, bottom=2.2cm, left=2.5cm, right=2.5cm}

% ─── Colores ─────────────────────────────────────────────────────────────────
\definecolor{azulmarco}{RGB}{30, 90, 160}
\definecolor{grisclaro}{RGB}{245, 246, 248}
\definecolor{naranjasupp}{RGB}{160, 60, 0}

% ─── Encabezados de sección ──────────────────────────────────────────────────
% comando auxiliar: titlesec le pasa el título como argumento (#1 es legal aquí)
\newcommand{\barrasec}[1]{%
  \colorbox{azulmarco}{\makebox[\dimexpr\linewidth-2\fboxsep\relax][l]{\quad\thesection\;\; #1}}}
\titleformat{\section}[block]
  {\normalfont\large\bfseries\color{white}}
  {}
  {0pt}
  {\barrasec}

\titleformat{\subsection}[hang]
  {\normalfont\bfseries\color{azulmarco}}
  {\thesubsection}
  {0.5em}
  {}

% Barra de sección sin numeración (para secciones de cierre como el resumen)
\newcommand{\barrasimple}[1]{%
  \par\bigskip\noindent
  \colorbox{azulmarco}{\makebox[\dimexpr\linewidth-2\fboxsep\relax][l]{%
    \normalfont\large\bfseries\color{white}\quad #1}}\par\medskip}

% ─── Entornos ────────────────────────────────────────────────────────────────
\newmdenv[
  backgroundcolor=grisclaro,
  linecolor=azulmarco,
  linewidth=1pt,
  innerleftmargin=10pt,
  innerrightmargin=10pt,
  innertopmargin=8pt,
  innerbottommargin=8pt,
  skipabove=4pt,
  skipbelow=4pt
]{demobox}

% Nota expositiva (comentario al margen del paso)
\newcommand{\nota}[1]{%
  \smallskip\par
  \noindent{\small\textit{Nota:} #1.}\par
  \smallskip
}

% Supuesto del modelo (se resalta en naranja cuando un paso deja de ser identidad)
\newcommand{\supuesto}[1]{%
  \smallskip\par
  \noindent{\small\textbf{\textcolor{naranjasupp}{Supuesto:}} \textit{#1.}}\par
  \smallskip
}

% Indicadora
\newcommand{\ind}[1]{\mathbf{1}\!\left[#1\right]}

% ─── Metadata ────────────────────────────────────────────────────────────────
\title{%
  \textbf{Construcci\'on del marco te\'orico 3.2}\\[4pt]
  \large Adopci\'on heterog\'enea de IA generativa\\con exposici\'on ling\"u\'istica diferencial\\[2pt]
  \normalsize Derivaci\'on formal paso a paso
}
\author{}
\date{}

% ─────────────────────────────────────────────────────────────────────────────
\begin{document}
\maketitle
\thispagestyle{empty}

\noindent
Este documento presenta, paso a paso, c\'omo se construye el marco te\'orico que
conduce a la ecuaci\'on estimable del modelo. Se parte de la decisi\'on de un
desarrollador a nivel de tarea y se agrega progresivamente hasta la
especificaci\'on de panel. A lo largo de la derivaci\'on se indica con claridad
en qu\'e momento se introduce cada \textbf{supuesto del modelo}, de modo que la
construcci\'on quede completamente transparente.

\bigskip

% ─── Notacion basica ─────────────────────────────────────────────────────────
\subsection*{Notaci\'on b\'asica}

\begin{center}
\renewcommand{\arraystretch}{1.2}
\begin{tabular}{ll}
\toprule
\textbf{S\'imbolo} & \textbf{Significado} \\
\midrule
$z \in [0,1]$ & Tarea elemental de programaci\'on \\
$h_c(z) > 0$ & Productividad humana en la tarea $z$ del lenguaje $c$ \\
$a_c(z) \ge 0$ & Productividad de ChatGPT en la tarea $z$ del lenguaje $c$ \\
$\mu_j \in \{0,1\}$ & Indicador de acceso del desarrollador $j$ a ChatGPT \\
$y_{jc}(\mu_j)$ & Producci\'on total del desarrollador $j$ en lenguaje $c$ \\
$\mathcal{Z}^{\mathrm{auto}}_c$ & Conjunto de tareas donde ChatGPT supera al humano \\
$\tau_c$ & Fracci\'on de tareas automatizables \\
$\rho_c$ & Ahorro relativo promedio en tareas automatizables \\
$\phi_i$ & Capacidad de absorci\'on tecnol\'ogica del pa\'is $i$ \\
$A_{it}$ & Disponibilidad de ChatGPT en pa\'is $i$, trimestre $t$ \\
$Y_{ict}$ & \textit{Unique pushers} por 100\,mil hab.\ en pa\'is $i$, lenguaje $c$, trimestre $t$ \\
\bottomrule
\end{tabular}
\end{center}

\bigskip

% =============================================================================
\section{Funci\'on de producci\'on de c\'odigo a nivel de tarea}

% -----------------------------------------------------------------------------
\subsection{Productividad efectiva por tarea}

Con acceso potencial a ChatGPT, el desarrollador elige el m\'etodo m\'as productivo:

\begin{demobox}
\begin{equation}
  y_{jc}(z;\,\mu_j) = \max\!\bigl\{h_c(z),\; \mu_j\, a_c(z)\bigr\}.
\end{equation}

\textbf{Si $\mu_j = 0$} (sin acceso a ChatGPT):
\begin{equation}
  y_{jc}(z;\,0) = \max\!\bigl\{h_c(z),\;0\bigr\} = h_c(z),
\end{equation}
porque $h_c(z) > 0$ por hip\'otesis.

\textbf{Si $\mu_j = 1$} (con acceso a ChatGPT):
\begin{equation}
  y_{jc}(z;\,1) = \max\!\bigl\{h_c(z),\;a_c(z)\bigr\}.
\end{equation}
\end{demobox}

% -----------------------------------------------------------------------------
\subsection{Producci\'on total}

La producci\'on total integra la productividad sobre el continuo de tareas $z \in [0,1]$:

\begin{demobox}
\begin{equation}
  y_{jc}(\mu_j) = \int_0^1 \max\!\bigl\{h_c(z),\;\mu_j\,a_c(z)\bigr\}\,dz.
\end{equation}

En los dos casos extremos:
\begin{align}
  y_{jc}(0) &= \int_0^1 h_c(z)\,dz, \\
  y_{jc}(1) &= \int_0^1 \max\!\bigl\{h_c(z),\;a_c(z)\bigr\}\,dz.
\end{align}

Adem\'as $y_{jc}(1) \ge y_{jc}(0)$ porque $\max\{h_c,a_c\} \ge h_c$ puntualmente: el
acceso a la IA nunca reduce la producci\'on.
\end{demobox}

% =============================================================================
\section{Conjunto de tareas automatizables}

% -----------------------------------------------------------------------------
\subsection{Definici\'on del conjunto}

\begin{demobox}
\begin{equation}
  \mathcal{Z}^{\mathrm{auto}}_c = \bigl\{z \in [0,1] : a_c(z) > h_c(z)\bigr\}.
\end{equation}
Contiene exactamente las tareas donde ChatGPT genera una \emph{ganancia marginal estricta}.
\end{demobox}

% -----------------------------------------------------------------------------
\subsection{Reescritura del m\'aximo con funci\'on indicadora}

La funci\'on indicadora $\ind{\cdot} \in \{0,1\}$ permite separar los dos reg\'imenes y eliminar el operador $\max$:

\begin{demobox}
\begin{equation}
  \max\!\bigl\{h_c(z),\,a_c(z)\bigr\}
  = h_c(z)\,\ind{z \notin \mathcal{Z}^{\mathrm{auto}}_c}
  + a_c(z)\,\ind{z \in \mathcal{Z}^{\mathrm{auto}}_c}.
\end{equation}

\textbf{Por casos:}
\begin{itemize}[leftmargin=1.4em, itemsep=2pt]
  \item Si $z \in \mathcal{Z}^{\mathrm{auto}}_c$: $a_c(z)>h_c(z)$, el m\'aximo es $a_c(z)$;
        las indicadoras valen $1$ y $0$ respectivamente. \checkmark
  \item Si $z \notin \mathcal{Z}^{\mathrm{auto}}_c$: $a_c(z)\le h_c(z)$, el m\'aximo es $h_c(z)$;
        las indicadoras valen $0$ y $1$ respectivamente. \checkmark
\end{itemize}
\end{demobox}

% -----------------------------------------------------------------------------
\subsection{Medida de tareas automatizables}

\begin{demobox}
\begin{equation}
  \tau_c = \int_0^1 \ind{z \in \mathcal{Z}^{\mathrm{auto}}_c}\,dz
         = \int_0^1 \ind{a_c(z)>h_c(z)}\,dz.
\end{equation}
Como $\ind{\cdot} \in \{0,1\}$ y el intervalo tiene longitud~$1$:
\[
  0 \le \tau_c \le 1.
\]
$\tau_c$ es, entonces, la \emph{fracci\'on} del trabajo que la IA puede automatizar.
\end{demobox}

% =============================================================================
\section{Ahorro relativo por tarea automatizable}

\begin{demobox}
Para $z \in \mathcal{Z}^{\mathrm{auto}}_c$ se define el \emph{ahorro relativo} que aporta la IA:
\begin{equation}
  s_c(z) = \frac{a_c(z) - h_c(z)}{h_c(z)}.
\end{equation}
Como $a_c(z)>h_c(z)>0$, se tiene $s_c(z)>0$.

El ahorro promedio \emph{condicional} a las tareas automatizables es:
\begin{equation}
  \rho_c = \frac{1}{\tau_c}\int_{\mathcal{Z}^{\mathrm{auto}}_c} s_c(z)\,dz, \qquad \tau_c > 0.
\end{equation}
\end{demobox}

\nota{si $\tau_c = 0$ no hay tareas automatizables y el efecto del lenguaje se interpreta como nulo}

% =============================================================================
\section{Ganancia individual de productividad con acceso a IA}

% -----------------------------------------------------------------------------
\subsection{Producci\'on con IA --- usando la partici\'on del Paso 2}

\begin{demobox}
Sustituyendo la reescritura con indicadoras en $y_{jc}(1)$:
\begin{align}
  y_{jc}(1) &= \int_0^1\!\Bigl[h_c(z)\,\ind{z \notin \mathcal{Z}^{\mathrm{auto}}_c}
               + a_c(z)\,\ind{z \in \mathcal{Z}^{\mathrm{auto}}_c}\Bigr]dz \\[4pt]
             &= \int_{[0,1]\setminus\mathcal{Z}^{\mathrm{auto}}_c}\!h_c(z)\,dz
               + \int_{\mathcal{Z}^{\mathrm{auto}}_c}\!a_c(z)\,dz.
\end{align}
Se usan: linealidad de la integral y la identidad
$\int_0^1 f\,\ind{z\in A}\,dz = \int_A f\,dz$.
\end{demobox}

% -----------------------------------------------------------------------------
\subsection{Producci\'on sin IA --- misma partici\'on}

\begin{demobox}
El intervalo $[0,1]$ se parte en tareas no automatizables y automatizables (conjuntos disjuntos):
\[
  [0,1] = \bigl([0,1]\setminus\mathcal{Z}^{\mathrm{auto}}_c\bigr)\cup\mathcal{Z}^{\mathrm{auto}}_c.
\]
\begin{equation}
  y_{jc}(0) = \int_{[0,1]\setminus\mathcal{Z}^{\mathrm{auto}}_c}\!h_c(z)\,dz
             + \int_{\mathcal{Z}^{\mathrm{auto}}_c}\!h_c(z)\,dz.
\end{equation}
\end{demobox}

% -----------------------------------------------------------------------------
\subsection{Ganancia individual exacta}

\begin{demobox}
\begin{align}
  \Delta y_{jc} &= y_{jc}(1) - y_{jc}(0) \notag\\
  &= \left[\int_{\mathcal{Z}^{\mathrm{auto}}_c}\!a_c(z)\,dz\right]
    - \left[\int_{\mathcal{Z}^{\mathrm{auto}}_c}\!h_c(z)\,dz\right] \\[4pt]
  &= \int_{\mathcal{Z}^{\mathrm{auto}}_c}\!\bigl[a_c(z)-h_c(z)\bigr]\,dz.
\end{align}
Las integrales sobre las tareas \emph{no automatizables} se cancelan al restar.
\end{demobox}

% -----------------------------------------------------------------------------
\subsection{Reescritura con el ahorro relativo}

\begin{demobox}
De $s_c(z)=\frac{a_c(z)-h_c(z)}{h_c(z)}$ se obtiene $a_c(z)-h_c(z)=h_c(z)\,s_c(z)$, luego:
\begin{equation}
  \Delta y_{jc} = \int_{\mathcal{Z}^{\mathrm{auto}}_c}\!h_c(z)\,s_c(z)\,dz.
\end{equation}
\end{demobox}

\nota{hasta este punto la construcci\'on es una cadena de identidades exactas; a\'un no se ha introducido ning\'un supuesto}

% -----------------------------------------------------------------------------
\subsection{Del producto de integrales al producto de promedios}

\begin{demobox}
Definamos los promedios condicionales sobre $\mathcal{Z}^{\mathrm{auto}}_c$:
\begin{equation}
  \bar{h}_c = \frac{1}{\tau_c}\int_{\mathcal{Z}^{\mathrm{auto}}_c}\!h_c(z)\,dz,
  \qquad
  \rho_c = \frac{1}{\tau_c}\int_{\mathcal{Z}^{\mathrm{auto}}_c}\!s_c(z)\,dz.
\end{equation}

Por la identidad de covarianza:
\begin{equation}
  \mathbb{E}_{\mathcal{Z}^{\mathrm{auto}}_c}[h_c\,s_c]
  = \mathbb{E}_{\mathcal{Z}^{\mathrm{auto}}_c}[h_c]\cdot\mathbb{E}_{\mathcal{Z}^{\mathrm{auto}}_c}[s_c]
  + \mathrm{Cov}_{\mathcal{Z}^{\mathrm{auto}}_c}(h_c,s_c),
\end{equation}
la \textbf{forma general} de la ganancia individual es:
\begin{equation}\label{eq:exacta}
  \Delta y_{jc} = \tau_c\!\left[\bar{h}_c\,\rho_c
  + \mathrm{Cov}_{\mathcal{Z}^{\mathrm{auto}}_c}(h_c,s_c)\right].
\end{equation}

El modelo adopta la forma multiplicativa $\Delta y_{jc} = \tau_c\,\rho_c\,\bar{h}_c$,
que corresponde al caso particular
\begin{equation}
  \mathrm{Cov}_{\mathcal{Z}^{\mathrm{auto}}_c}(h_c,s_c) = 0.
\end{equation}
\end{demobox}

\supuesto{productividad humana $h_c(z)$ y ahorro relativo $s_c(z)$ no correlacionados dentro de $\mathcal{Z}^{\mathrm{auto}}_c$. Es el primer supuesto del modelo (la ecuaci\'on~\eqref{eq:exacta} es la forma general que lo contiene como caso particular)}

% -----------------------------------------------------------------------------
\subsection{Normalizaci\'on de escala}

\begin{demobox}
Se elige la unidad de producci\'on tal que el promedio de productividad humana en
las tareas automatizables sea unitario, $\bar{h}_c = 1$:
\begin{equation}
  \Delta y_{jc} = \tau_c\,\rho_c.
\end{equation}
La ganancia individual queda como el producto \textbf{fracci\'on automatizable}
($\tau_c$) $\times$ \textbf{ahorro relativo promedio} ($\rho_c$).
\end{demobox}

\nota{$\bar{h}_c = 1$ no es un resultado emp\'irico, es una elecci\'on de escala}

% =============================================================================
\section{Decisi\'on de entrada al mercado de programadores}

\begin{demobox}
La utilidad de programar es $U_{\mathrm{prog}}(\eta;\mu_j) = \eta + y_{jc}(\mu_j)$
y la utilidad de reserva es $U_{\mathrm{out}} = r_i$.

El individuo entra si $\eta + y_{jc}(\mu_j) \ge r_i$.
Definiendo los \emph{umbrales m\'inimos de habilidad}:
\begin{equation}
  \eta^*_c(0) = r_i - y_{jc}(0), \qquad
  \eta^*_c(1) = r_i - y_{jc}(1).
\end{equation}

La diferencia de umbrales es:
\begin{equation}
  \eta^*_c(0) - \eta^*_c(1)
  = \bigl[r_i - y_{jc}(0)\bigr] - \bigl[r_i - y_{jc}(1)\bigr]
  = \Delta y_{jc} = \tau_c\rho_c.
\end{equation}

Por tanto $\eta^*_c(1) = \eta^*_c(0) - \tau_c\rho_c$: el acceso a la IA \textbf{reduce}
el umbral de habilidad m\'inima, de modo que entran programadores marginales.
\medskip

\textit{Nota:} si se incluye la covarianza exacta del Paso~4.5,
$\eta^*_c(1) = \eta^*_c(0) - \tau_c\!\bigl[\bar{h}_c\rho_c + \mathrm{Cov}(h_c,s_c)\bigr]$.
\end{demobox}

% =============================================================================
\section{Agregaci\'on a desarrolladores activos}

\begin{demobox}
Sea $N_i$ la poblaci\'on del pa\'is $i$ y $F_i$ la CDF de habilidad $\eta$.
\begin{equation}
  \mathrm{Devs}_{ic}(\mu) = N_i\,\Pr[\eta \ge \eta^*_c(\mu)]
  = N_i\bigl[1-F_i(\eta^*_c(\mu))\bigr].
\end{equation}

La diferencia entre r\'egimenes es:
\begin{equation}
  \mathrm{Devs}_{ic}(1) - \mathrm{Devs}_{ic}(0)
  = N_i\bigl[F_i(\eta^*_c(0)) - F_i(\eta^*_c(1))\bigr].
\end{equation}

Sea $\Delta\eta = \eta^*_c(0)-\eta^*_c(1)=\tau_c\rho_c$.
Aplicando Taylor de primer orden alrededor de $\eta^*_c(0)$:
\begin{equation}
  F_i(\eta^*_c(0)-\Delta\eta) \approx F_i(\eta^*_c(0)) - f_i(\eta^*_c(0))\,\Delta\eta,
\end{equation}
se obtiene:
\begin{equation}
  \mathrm{Devs}_{ic}(1) - \mathrm{Devs}_{ic}(0)
  \approx N_i\,f_i(\eta^*_c(0))\,\tau_c\rho_c.
\end{equation}
\end{demobox}

\supuesto{la aproximaci\'on de Taylor de primer orden es v\'alida si $\Delta\eta = \tau_c\rho_c$ es peque\~no relativo a la escala en que cambia la densidad $f_i$}

% =============================================================================
\section{Conversi\'on per c\'apita y capacidad de absorci\'on}

\begin{demobox}
Definimos el resultado por cada 100\,000 habitantes:
\begin{equation}
  Y_{ic} = \frac{\mathrm{Devs}_{ic}}{N_i}\cdot 10^5.
\end{equation}

Entonces:
\begin{equation}
  \Delta Y_{ic}
  = \frac{\mathrm{Devs}_{ic}(1)-\mathrm{Devs}_{ic}(0)}{N_i}\cdot 10^5
  \approx 10^5\,f_i(\eta^*_c(0))\,\tau_c\rho_c.
\end{equation}

Se define la \textbf{capacidad de absorci\'on tecnol\'ogica} del pa\'is $i$:
\begin{equation}
  \phi_i = 10^5\,f_i(\eta^*_c(0)),
\end{equation}
y la \textbf{exposici\'on del lenguaje} $c$:
\begin{equation}
  \theta_c = \tau_c\rho_c.
\end{equation}

Por tanto: $\Delta Y_{ic} \approx \theta_c\,\phi_i$.
\end{demobox}

% -----------------------------------------------------------------------------
\subsection{Separabilidad pa\'is--lenguaje en $\phi_i$}

\begin{demobox}
El umbral $\eta^*_c(0) = r_i - y_{jc}(0)$ \textbf{depende del lenguaje} $c$, ya que
$y_{jc}(0)=\int_0^1 h_c(z)\,dz$ var\'ia con $c$. En rigor, entonces, la densidad
evaluada en el umbral es espec\'ifica del par $(i,c)$:
\begin{equation}
  \phi_{ic} = 10^5\,f_i\!\bigl(\eta^*_c(0)\bigr).
\end{equation}

Escribir $\phi_i$ (solo pa\'is) y obtener la forma multiplicativa limpia
$\Delta Y_{ic} = \theta_c\,\phi_i$ requiere que
\begin{equation}
  f_i\!\bigl(\eta^*_c(0)\bigr) \approx f_i \quad (\text{constante en } c),
\end{equation}
es decir, que la densidad de habilidades sea aproximadamente plana en el rango de
umbrales que inducen los distintos lenguajes (o, de modo equivalente, que las
productividades base $y_{jc}(0)$ sean similares entre lenguajes).
\end{demobox}

\supuesto{separabilidad pa\'is--lenguaje. Sin ella la capacidad de absorci\'on ser\'ia $\phi_{ic}$ y la forma $\theta_c\phi_i$ no se separar\'ia limpiamente; el ATT por lenguaje sigue identificado, pero como $\theta_c\,\bar{\phi}_{\mathcal{T},c}$ (promedio de absorci\'on espec\'ifico del lenguaje)}

% =============================================================================
\section{Especificaci\'on de panel con efectos fijos}

\begin{demobox}
El indicador de disponibilidad es $A_{it} \in \{0,1\}$.

La l\'inea base sin tratamiento se descompone como:
\begin{equation}
  Y_{ic}(0) = \alpha_i + \beta_t + \delta_c + \eta_{ict},
\end{equation}
donde $\alpha_i$ son efectos pa\'is, $\beta_t$ efectos tiempo, $\delta_c$ efectos lenguaje.

Sustituyendo $\Delta Y_{ic} \approx \theta_c\phi_i\,A_{it}$:
\begin{equation}
  \boxed{Y_{ict} = \alpha_i + \beta_t + \delta_c + \theta_c\,\phi_i\,A_{it} + \eta_{ict}.}
\end{equation}

Para interpretaci\'on \emph{causal} se requiere exogeneidad condicional:
\[
  \mathbb{E}[\eta_{ict}\mid\alpha_i,\beta_t,\delta_c,A_{it}] = 0.
\]
\end{demobox}

% =============================================================================
\barrasimple{Resumen paso a paso}

\begin{center}
\renewcommand{\arraystretch}{1.4}
\begin{tabular}{>{\bfseries}p{1.6cm} p{6.2cm} p{4.8cm}}
\toprule
Paso & Resultado & Tipo \\
\midrule
1 & $y_{jc}(\mu_j)=\int_0^1\max\{h_c(z),\mu_j a_c(z)\}dz$ & Identidad / definici\'on \\
2 & $\mathcal{Z}^{\mathrm{auto}}_c=\{z:a_c(z)>h_c(z)\}$;\; $\tau_c=\int_0^1\ind{z\in\mathcal{Z}^{\mathrm{auto}}_c}dz$ & Definici\'on exacta \\
3 & $s_c(z)=(a_c-h_c)/h_c$;\; $\rho_c=\frac{1}{\tau_c}\int s_c\,dz$ & Definici\'on exacta ($\tau_c>0$)\\
4.1--4.4 & $\Delta y_{jc}=\int_{\mathcal{Z}^{\mathrm{auto}}_c}h_c s_c\,dz$ & Identidad exacta \\
4.5 & $\Delta y_{jc}=\tau_c\rho_c\bar{h}_c$ & \textcolor{naranjasupp}{\textbf{Supuesto}: $\mathrm{Cov}(h_c,s_c)=0$} \\
4.6 & $\Delta y_{jc}=\tau_c\rho_c$ & Normalizaci\'on $\bar{h}_c=1$ \\
5 & $\eta^*_c(1)=\eta^*_c(0)-\tau_c\rho_c$ & Depende del Paso 4 \\
6 & $\mathrm{Devs}_{ic}(1)-\mathrm{Devs}_{ic}(0)\approx N_i f_i(\eta^*_c(0))\tau_c\rho_c$ & Taylor 1.\textordmasculine\ orden \\
7 & $\Delta Y_{ic}\approx\tau_c\rho_c\phi_i$ & Aprox.\ + \textcolor{naranjasupp}{\textbf{separabilidad} $\phi_{ic}\!\approx\!\phi_i$} \\
8 & $Y_{ict}=\alpha_i+\beta_t+\delta_c+\theta_c\phi_i A_{it}+\eta_{ict}$ & Especificaci\'on de panel \\
\bottomrule
\end{tabular}
\end{center}

\smallskip
\noindent{\small En \textcolor{naranjasupp}{naranja}, los pasos en que se introduce un supuesto del modelo; el resto son identidades o definiciones exactas.}

\bigskip

\subsection*{S\'intesis de la derivaci\'on}

La construcci\'on no requiere ning\'un supuesto hasta
\[
  \Delta y_{jc} = \int_{\mathcal{Z}^{\mathrm{auto}}_c}\!h_c(z)\,s_c(z)\,dz.
\]

De aqu\'i en adelante se introducen los supuestos del modelo. La forma general de la
ganancia individual es
\[
  \Delta y_{jc} = \tau_c\!\left[\bar{h}_c\,\rho_c + \mathrm{Cov}_{\mathcal{Z}^{\mathrm{auto}}_c}(h_c,s_c)\right],
\]
de la cual $\tau_c\rho_c\bar{h}_c$ es el caso particular con covarianza nula. La ecuaci\'on
estimable final,
\[
  Y_{ict} = \alpha_i + \beta_t + \delta_c + \theta_c\,\phi_i\,A_{it} + \eta_{ict},
\]
se obtiene como modelo estructural simplificado a partir de \textbf{cuatro} supuestos:
\begin{enumerate}[leftmargin=2em, itemsep=2pt]
  \item \textbf{Covarianza cero} entre productividad humana $h_c(z)$ y ahorro relativo $s_c(z)$
        dentro de las tareas automatizables.
  \item \textbf{Normalizaci\'on} $\bar{h}_c = 1$ (elecci\'on de escala).
  \item \textbf{Aproximaci\'on de Taylor de primer orden} para convertir la reducci\'on del umbral
        en aumento de desarrolladores activos.
  \item \textbf{Separabilidad pa\'is--lenguaje}: que la densidad en el umbral $f_i(\eta^*_c(0))$
        sea aproximadamente independiente del lenguaje $c$, de modo que $\phi_{ic}\approx\phi_i$
        y el coeficiente del tratamiento se factorice limpiamente como $\theta_c\,\phi_i$.
\end{enumerate}

\end{document}
